\chapter{Lecture 03}
\hfill {\footnotesize Date: 15/02/2026}
\section{Prime \& Maximal Ideals}
\begin{definition}[Prime Ideal]
    An ideal \(\mathfrak{p}\) in \(R\) is called \emph{prime} if \(\mathfrak{p} \neq (1)\) and if \(xy \in \mathfrak{p}\) implies \(x \in \mathfrak{p}\)  or \(y \in \mathfrak{p}\).   
\end{definition}
\begin{example}
    Take the ring \(R := (\mathbb{Z}, +, \times)\). Then prime ideals of \(R\) are precisely the ideals \((p)\), where \(p\) is prime.

    Because, first of all, every ideal of \(R\) is of the form \((n)\). 

    And we claim that ``\((n)\) is prime ideal iff \(n\) is a prime number;'' 
    equivalently, ``\(\forall a, b \in \mathbb{Z}, ab \in (n) \implies a \in (n) \text{ or } b \in (n)\) iff \(n\) is a prime number;''
 equivalently, ``\(\forall a, b \in Z, n \vert ab \implies n \vert a \text{ or } n \vert b\) iff \(n\) is a prime number.'' We prove that the last statement is indeed true.
    \begin{claim}
        For all \(a, b \in \mathbb{Z}, n \vert ab \implies n \vert a \text{ or } n \vert b\) iff \(n\) is a prime number. 
    \end{claim}
    \begin{proof}[Proof of claim]
        Let for all \(a, b \in \mathbb{Z}, n | ab \implies n | a \text{ or } n | b\). And suppose \(n\) is composite. Then there exists integers \(a'\neq 1, b' \neq 1\) such that \(n = a'b'\). Then \(n = a'b' \big| a' \frac{\text{ lcm}(a', b')}{a'}\) but \(a'b' \not| a'\) and \(a'b' \not| \frac{\text{ lcm}(a', b')}{a'}\), contradicting the hypothesis. Hence, \(n\) must be prime.

        Conversely, let \(n\) be prime. If \(n | a\), we are done. Otherwise, \(n \not| a \implies \operatorname{gcd}(n,a) = 1\), which implies \(px + ay = 1 \implies pbx + aby = b\) for some \(x, y \in \mathbb{Z}\)\footnote{Bezout's identity.}. Since \(p | pbx\) and \(p | (ab)y\)\footnote{since \(p | ab\) by hypothesis.}, therefore \(p | b\).
    \end{proof}
    Therefore, prime ideals in \(\mathbb{Z}\) are principle ideals generated by prime numbers.
\end{example}
\begin{definition}[Maximal Ideal]
    An ideal \(\mathfrak{m}\) in \(R\) is  called \emph{maximal} if \(\mathfrak{m} \neq (1)\) and if there is no ideal \(\mathfrak{a}\) in \(R\) such that \[\mathfrak{m} \subset \mathfrak{a} \subset R \quad \text{(strict inclusion).}\]
    In other words, the only ideal that strictly contains \(\mathfrak{m}\) is \((1)\).
\end{definition}
\begin{proposition}
    \begin{enumerate}
        \item An ideal \(\mathfrak{p}\) is prime if and only if \(R/\mathfrak{p}\) is an integral domain.
        \item An ideal \(\mathfrak{m}\) is maximal if and only if \(R/\mathfrak{m}\) is a field.
    \end{enumerate}
\end{proposition}
\begin{proof}
    \begin{enumerate}
        \item Assume \(\mathfrak{p}\) is a prime ideal. Let \(a+\mathfrak{p}, b+\mathfrak{p} \in R/\mathfrak{p}\) and \((a + \mathfrak{p}) \cdot (b+\mathfrak{p}) = ab + \mathfrak{p}= 0\footnote{It is clear from the context that \(0 \in R/\mathfrak{p}\).}\). \( ab + \mathfrak{p} = 0\) implies \(ab \in \mathfrak{p}\). Since \(\mathfrak{p}\) is prime, \(a \in \mathfrak{p} \implies a+\mathfrak{p} = 0\) or \(b \in \mathfrak{p} \implies b + \mathfrak{p} = 0\). So, only zero-divisor in \(R/\mathfrak{p}\) is \(0\), i.e., \(R/\mathfrak{p}\) is an integral domain.

        Conversely, assume that \(R/\mathfrak{p}\) is an integral domain. Let \(ab \in \mathfrak{p}\). Then \(0 = ab + \mathfrak{p} = (a + \mathfrak{p} )(b + \mathfrak{p})\), which implies that \(a + \mathfrak{p} = 0 \) or \(b + \mathfrak{p} = 0 \), since \(R/\mathfrak{p}\) is an integral domain. Thus, \(a \in \mathfrak{p}\) or \(b \in \mathfrak{p}\) and so \(\mathfrak{p}\) is a prime ideal.

    \item Suppose \(\mathfrak{m}\) is a maximal ideal of \(R\).
        By the \hyperref[proposition:correspondencethm]{Correspondence Theorem}, there is a bijection between the set of ideals 
\(\mathfrak{a}\) of \(R\) containing \(\mathfrak{m}\) and the set of ideals of 
\(R/\mathfrak{m}\).
Since \(\mathfrak{m}\) is maximal, the only ideals of \(R\) containing 
\(\mathfrak{m}\) are: \(\mathfrak{m}\) and \(R\). Therefore, the only ideals 
of \(R/\mathfrak{m}\) are: \((0), (1)\).
Hence, \(R/\mathfrak{m}\) has no nontrivial ideals, so it is a field.

 Conversely, suppose \(R/\mathfrak{m}\) is a field.
Then its only ideals are \((0)\) and \(R/\mathfrak{m}\). By the \hyperref[proposition:correspondencethm]{Correspondence Theorem}, the ideals of \(R/\mathfrak{m}\) correspond bijectively to the ideals 
of \(R\) containing \(\mathfrak{m}\). Therefore, the only ideals of \(R\) 
containing \(\mathfrak{m}\) are: \(\mathfrak{m}\) and \(R\). Hence, 
\(\mathfrak{m}\) is a maximal ideal of \(R\).
    \end{enumerate}
\end{proof}
\begin{corollary}
    Every maximal ideal is prime, but the converse does not hold true.
\end{corollary}
\begin{proof}
    Let \(\mathfrak{m}\) be a maximal ideal in \(R\). Then \(R/\mathfrak{m}\) is a field. Then \(R/\mathfrak{m}\) is an integral domain automatically. Therefore, \(\mathfrak{m}\) is a prime ideal.

    But \((0)\) is a prime ideal in \(\mathbb{Z}\), but not maximal ideal because \((0)\subset (p) \subset (1)\) (strict inclusion).
\end{proof}
\begin{theorem}
    Suppose \(f: R \longrightarrow R'\) is a ring homomorphism, and let \(\mathfrak{p}'\) be a prime ideal in \(R'\). Then, \(f^{-1}(\mathfrak{p}')\) -- the preimage of the ideal is a prime ideal in \(R\).
\end{theorem}
\begin{proof}
    Suppose \(xy \in f^{-1}(\mathfrak{p}')\), which implies \(f(xy) \in \mathfrak{p}'\), and that implies \(f(x)f(y) \in \mathfrak{p}'\), and that implies \(f(x) \in \mathfrak{p}' \text{ or } f(y) \in \mathfrak{p}'\), and that implies \(x \in f^{-1}(\mathfrak{p}') \text{ or } y \in f^{-1}(y)\), which proves \(f^{-1}(\mathfrak{p}')\) is a prime ideal.
\end{proof}
But preimage of a maximal ideal may not be maximal.
\begin{proof}[Counter-example]
    Consider the inclusion map \(\iota: \mathbb{Z} \longrightarrow \mathbb{Q}\). Since \(\mathbb{Q}\) is a field, its only ideals are: \((0)\) and \((1)\). So \((0)\) is a maximal ideal. But \(\iota^{-1}((0)) = (0) \subset (2) \subset (1)\) is not maximal in \(\mathbb{Z}\). 
\end{proof}
\begin{theorem}[Existence of a Maximal Ideal]\label{theorem:maximalidealexistence}
    Every ring \(R \neq (0)\) has at least one maximal ideal.
\end{theorem}
\begin{proof}
    We use Zorn's lemma.
    \begin{theorem}[Zorn's Lemma]
    Let \((S, \le)\) be a non-empty partially ordered set
        \footnote{which means that \(\le\) is a relation on \(S\) satisfying:               \begin{itemize}
                \item \(x \le x\) for all \(x \in S\), 
                \item \(x \le y\) and \(y \le x\) implies \(x = y\) for all \(x, y \in S\),
                \item \(x \le y, y \le z \implies x \le z\) for all \(x, y, z \in S\).
            \end{itemize}
        } (also known as poset). 
        If every chain\footnote{A \textit{chain} in \(S\) is a subset \(C\) of \(S\) such that for any \(x, y \in C\) we have either \(x \le y\) or \(y \le x\).} in \(S\) has an upper bound\footnote{an element \(u\) such that \(x \le u\) for all \(x \in C\)} in \(S\), then \(S\) has at least one maximal element.
    \end{theorem}
    \begin{proof}[Proof of Zorn's lemma]Omitted.\end{proof}
    Let \(\Sigma\) be the set of all ideals \(\mathfrak{a} \neq (1)\) and define partial order \(\subseteq\) using set inclusion. So \((\Sigma, \subseteq)\) is a poset. To apply Zorn's lemma, we need to show that every chain in \((\Sigma, \subseteq)\) has an upper bound in \(\Sigma\).

    Suppose \(\mathcal{C}:= \{\mathfrak{a}_\alpha\}_{\alpha \in J}\) is an arbitrary chain. So \(\mathfrak{a}_\alpha \subseteq \mathfrak{a}_\beta\) or \(\mathfrak{a}_\beta \subseteq \mathfrak{a}_\alpha\) for all \(\alpha, \beta \in J\). Define \[\mathfrak{a} := \bigcup_{\alpha \in J} \mathfrak{a}_\alpha,\] which qualifies as an ideal because \(\mathfrak{a}\) is a subgroup of \(R\) and closed under multiplication by any elements of \(R\):
    \begin{itemize}
        \item \(0 \in \mathfrak{a}_\alpha\) for all \(\alpha \in J\), which implies \(0 \in \mathfrak{a}\).
        \item Suppose \(x, y \in \mathfrak{a}\). So \(x \in \mathfrak{a}_\alpha, y \in \mathfrak{a}_\beta\) for some \(\alpha, \beta \in J\). Since \(\mathfrak{a}_\alpha, \mathfrak{a}_\beta \in \mathcal{C}\), without loss of generality, let \(\mathfrak{a}_\alpha \subseteq \mathfrak{a}_\beta\). Then \(x - y \in \mathfrak{a}_\beta\), implying \(x - y \in \mathfrak{a}\).
        \item Suppose \(r \in R\) and suppose \(x \in \mathfrak{a}\). That means \(x \in \mathfrak{a}_\alpha\) for some \(\alpha \in J\). Hence, \(rx \in \mathfrak{a}_\alpha \subseteq \mathfrak{a}\). 
    \end{itemize}
    Clearly, \(\mathfrak{a}_\alpha \subset \mathfrak{a}\) and \(\mathfrak{a} \in \Sigma\), i.e., every chain in \(\Sigma\) has an upper bound in \(\Sigma\). Therefore, Zorn's lemma implies that \(\Sigma\) has a maximal element. Hence, there is at least one maximal ideal of \(R\).
\end{proof}
\begin{corollary}\label{corollary:primeidealexistence}
    Every ring \(R \neq (0)\) has at least one prime ideal as well.
\end{corollary}
\begin{corollary}\label{corollary:idealcontainedinmaximal}
    If \(\mathfrak{a} \neq (1)\) is an ideal of \(R \neq (0)\), then there exists a maximal ideal \(\mathfrak{m}\) containing \(\mathfrak{a}\).
\end{corollary}
\begin{proof}[Proof 1]
\autoref{theorem:maximalidealexistence} guarantees existence of a maximal ideal \(\mathfrak{m}_q\) in \(R/\mathfrak{a}\). Using \hyperref[proposition:correspondencethm]{Correspondence Theorem}, we have a maximal ideal \(\mathfrak{m}\) such that \[\mathfrak{a} \subseteq \mathfrak{m} \subset R \quad (\text{second inclusion is strict)}.\] 
Maximality of \( \mathfrak{m} \) is guaranteed because the bijection in the correspondence theorem is order-preserving.
\end{proof}
\begin{proof}[Proof 2]
    \todo{To be written...}
\end{proof}
\begin{corollary}\label{corollary:nonunitinmaximalideal}
    Every non-unit of a ring \(R\) is contained in a maximal ideal.
\end{corollary}
\begin{proof}
    Suppose \(x \in R\) is a non-unit. Then \((x) \neq (1)\) because if we had \((x) = (1)\), then there would exist \(y \in R\) such that \(xy = 1\), which contradicts the fact that \(x\) is a non-unit. Using \autoref{corollary:idealcontainedinmaximal}, we conclude that there exists a maximal ideal \(\mathfrak{m}\) such that \(x \in (x) \subset \mathfrak{m}\).
\end{proof}
\section{Local Rings}
\begin{definition}[Local Ring]
    A ring \(R\) is called \emph{local} if it has only one maximal ideal \(\mathfrak{m}\).
\end{definition}
\begin{example}
    \begin{itemize}
        \item All fields are local rings, because a field has only two ideals: \((0)\) and \((1)\), so only maximal is \((0)\).
        \item \(\mathbb{Z}_n\) (also denoted by \(\mathbb{Z}/n\mathbb{Z}\)) is a local ring if and only if \(n = p^a\) for some prime \(p\) and positive integer \(a\).
            \begin{proof}
                
            \end{proof}

%            \begin{proof}
%                It is clear that all the ideals in \(\mathbb{Z}_{p^a}\) are of the form \((n)\), for \(n \in \mathbb{Z}_{p^a}\).
%
%                An ideal \(\mathfrak{a} \subset \mathbb{Z}_{p^a}\) is not equal \((1)\) iff \(n \in \mathbb{Z}_{p^a}\) is a non-unit. It is basic result from number theory that \(u \in \mathbb{Z}_{n}\) is a unit iff \(u\) is co-prime to \(n\), in particular, \(u \in \mathbb{Z}_{p^a}\) is a non-unit iff \(n | p\). Non-units of \(\mathbb{Z}_{p^a}\) are (marked in pink):
%                \[0, 1, \cdots, \colorbox{pink}{\(p\)}, p+1, \cdots, 2p-1, 2p, 2p+1, \cdots, p^2-1, \colorbox{pink}{\(p^2\)}, p^2+1, \cdots, kp^\alpha-1, 
%                \colorbox{pink}{\(kp^\alpha\)}, kp^\alpha+1, \cdots, p^n-1\] for \(\alpha \in \{1, \cdots, n-1\}\) and \(k \in \{1, \cdots, p-1\}\).


%            \end{proof}
    \end{itemize}
\end{example}
\todo{To be written...}
\begin{nonexample}
    \(\mathbb{Z}\) is not a local ring because all prime ideals are maximal in \(\mathbb{Z}\), and there are more than \( 1 \) prime ideal in  \( \mathbb{Z} \).
\end{nonexample}
\begin{definition}[Semi-local Ring]
    A ring \(R\) is called \emph{semi-local} if number of maximal ideals in \(R\) is finite.
\end{definition}
\begin{example}
    \(Z_n\) is always semi-local since its number of ideals is finite, so is number of maximal ideals.
\end{example}
\begin{exercise}
    Is \(\mathbb{K}[x]\) local? When is \(\mathbb{K}[x]/(x^n)\) local?

\end{exercise}
